\documentclass{beamer}

% Theme choice
\usetheme{Madrid} % You can change to e.g., Warsaw, Berlin, CambridgeUS, etc.

% Encoding and font
\usepackage[utf8]{inputenc}
\usepackage[T1]{fontenc}

% Graphics and tables
\usepackage{graphicx}
\usepackage{booktabs}

% Code listings
\usepackage{listings}
\lstset{
basicstyle=\ttfamily\small,
keywordstyle=\color{blue},
commentstyle=\color{gray},
stringstyle=\color{red},
breaklines=true,
frame=single
}

% Math packages
\usepackage{amsmath}
\usepackage{amssymb}

% Colors
\usepackage{xcolor}

% TikZ and PGFPlots
\usepackage{tikz}
\usepackage{pgfplots}
\pgfplotsset{compat=1.18}
\usetikzlibrary{positioning}

% Hyperlinks
\usepackage{hyperref}

% Title information
\title{Chapter 6: Data Structures: Dictionaries and Sets}
\author{Your Name}
\institute{Your Institution}
\date{\today}

\begin{document}

\frame{\titlepage}

\begin{frame}[fragile]
    \frametitle{Introduction to Dictionaries and Sets}
    \begin{block}{Overview}
        Overview of dictionaries and sets in Python, their importance in data manipulation, and their applications in group projects.
    \end{block}
\end{frame}

\begin{frame}[fragile]
    \frametitle{What are Dictionaries?}
    \begin{itemize}
        \item \textbf{Definition}: A dictionary is a mutable, unordered collection of key-value pairs. Each key is unique and is used to access its corresponding value.
        \item \textbf{Syntax}: Defined using curly braces \{\} with keys and values separated by a colon \(:\).
    \end{itemize}
    \begin{block}{Example}
        \begin{lstlisting}[language=Python]
student_info = {
    'name': 'Alice',
    'age': 21,
    'major': 'Computer Science'
}
        \end{lstlisting}
    \end{block}
\end{frame}

\begin{frame}[fragile]
    \frametitle{Key Features of Dictionaries}
    \begin{itemize}
        \item \textbf{Mutable}: Can change, add, or remove items after it is created.
        \item \textbf{Unordered}: Items do not maintain any specific order.
        \item \textbf{Key-Value Association}: Each value is accessed using its unique key.
    \end{itemize}
\end{frame}

\begin{frame}[fragile]
    \frametitle{What are Sets?}
    \begin{itemize}
        \item \textbf{Definition}: A set is an unordered collection of unique items. Sets are commonly used to remove duplicates from a list or to perform mathematical set operations.
        \item \textbf{Syntax}: Defined using curly braces \{\} or the \texttt{set()} function.
    \end{itemize}
    \begin{block}{Example}
        \begin{lstlisting}[language=Python]
fruits = {'apple', 'banana', 'orange'}
        \end{lstlisting}
    \end{block}
\end{frame}

\begin{frame}[fragile]
    \frametitle{Key Features of Sets}
    \begin{itemize}
        \item \textbf{Unordered}: The elements have no specific order, making it efficient for membership testing.
        \item \textbf{Unique Elements}: Automatically ensures that all elements in the set are unique.
        \item \textbf{Mutable}: Elements can be added or removed from a set.
    \end{itemize}
\end{frame}

\begin{frame}[fragile]
    \frametitle{Importance in Data Manipulation}
    \begin{itemize}
        \item \textbf{Data Organization}: Both dictionaries and sets allow for efficient organization of data, making it easier to retrieve and manipulate.
        \item \textbf{Enhanced Performance}: Searching for elements in a set is faster (average O(1) time complexity) than in a list (average O(n) time complexity).
    \end{itemize}
\end{frame}

\begin{frame}[fragile]
    \frametitle{Applications in Group Projects}
    \begin{itemize}
        \item \textbf{Data Categorization}: Use dictionaries to store group members and their respective roles (e.g., \{'member1': 'researcher', 'member2': 'developer'\}).
        \item \textbf{Managing Unique Entries}: Utilize sets to keep track of completed tasks or unique contributors without duplicates.
    \end{itemize}
\end{frame}

\begin{frame}[fragile]
    \frametitle{Code Snippet Example}
    \begin{block}{Creating a Dictionary and Set}
        \begin{lstlisting}[language=Python]
# Creating a dictionary of group members
group_members = {
    'John': 'Project Manager',
    'Sara': 'Designer',
    'Mike': 'Developer'
}

# Creating a set of completed tasks
completed_tasks = {'design', 'development', 'testing'}

# Adding a new task to the set
completed_tasks.add('deployment')
        \end{lstlisting}
    \end{block}
\end{frame}

\begin{frame}[fragile]
    \frametitle{Takeaway}
    Mastering dictionaries and sets is essential for effective data manipulation in Python, aiding in logical structuring of group project tasks and roles.
\end{frame}

\begin{frame}[fragile]
    \frametitle{What is a Dictionary? - Definition}
    A \textbf{Dictionary} in Python is a built-in data structure that allows you to store data in a key-value pair format. This means that each value in the dictionary is associated with a unique key, which serves as an identifier.
    
    \begin{block}{Key Characteristics}
        \begin{itemize}
            \item \textbf{Key-Value Pairs}: Each entry consists of a key and its corresponding value, separated by a colon. 
            \item \textbf{Uniqueness of Keys}: Keys must be unique; assigning a new value to an existing key overrides the previous value.
            \item \textbf{Unordered}: Dictionaries do not maintain any order for the key-value pairs (prior to Python 3.7).
        \end{itemize}
    \end{block}
\end{frame}

\begin{frame}[fragile]
    \frametitle{What is a Dictionary? - Comparison with Lists}
    \begin{block}{How Dictionaries Differ from Lists}
        \begin{itemize}
            \item \textbf{Data Structure Type}:
                \begin{itemize}
                    \item \textbf{Lists}: Ordered collection, indexed by their position (e.g., \texttt{my\_list = [10, "Python", 3.14]}).
                    \item \textbf{Dictionaries}: Unordered collection, indexed by keys (e.g., \texttt{my\_dict = \{"fruit": "apple", "quantity": 5\}}).
                \end{itemize}
            \item \textbf{Access Method}:
                \begin{itemize}
                    \item \textbf{Lists}: Accessed by index (e.g., \texttt{my\_list[1]} returns ``Python'').
                    \item \textbf{Dictionaries}: Accessed by keys (e.g., \texttt{my\_dict["fruit"]} returns ``apple'').
                \end{itemize}
        \end{itemize}
    \end{block}
\end{frame}

\begin{frame}[fragile]
    \frametitle{Example of a Dictionary}
    \begin{block}{Example Dictionary}
        \begin{lstlisting}[language=Python]
student_info = {
    "name": "John",
    "age": 21,
    "major": "Computer Science",
    "gpa": 3.75
}
        \end{lstlisting}
    \end{block}
    
    \begin{itemize}
        \item \textbf{Keys}: ``name'', ``age'', ``major'', ``gpa''
        \item \textbf{Values}: ``John'', 21, ``Computer Science'', 3.75
    \end{itemize}
    
    \begin{block}{Key Points to Emphasize}
        \begin{itemize}
            \item \textbf{Fast Lookups}: Accessing a value by its key is efficient (O(1)).
            \item \textbf{Dynamic Size}: Dictionaries can grow and shrink as key-value pairs are added or removed.
            \item \textbf{Versatile Data}: Values can be of any data type; keys must be immutable (e.g., strings, numbers).
        \end{itemize}
    \end{block}
\end{frame}

\begin{frame}[fragile]
    \frametitle{Creating and Accessing Dictionaries - Overview}
    \begin{block}{Understanding Dictionaries}
        \begin{itemize}
            \item A dictionary in Python is a collection of key-value pairs.
            \item It is mutable, meaning you can change its contents after creation.
            \item Keys must be unique and immutable (strings, numbers, tuples) while values can be of any data type.
        \end{itemize}
    \end{block}
\end{frame}

\begin{frame}[fragile]
    \frametitle{Creating a Dictionary}
    \begin{block}{Methods of Creation}
        \begin{itemize}
            \item Use curly braces \{\} or the \texttt{dict()} constructor.
        \end{itemize}
    \end{block}

    \begin{block}{Example}
    \begin{lstlisting}[language=Python]
# Using curly braces
student = {
    "name": "Alice",
    "age": 22,
    "major": "Computer Science"
}

# Using dict() constructor
student = dict(name="Alice", age=22, major="Computer Science")
    \end{lstlisting}
    \end{block}
\end{frame}

\begin{frame}[fragile]
    \frametitle{Modifying Dictionaries}
    \begin{block}{Adding Items}
        Add a new key-value pair by assigning a value to a new key.
        \begin{lstlisting}[language=Python]
student["GPA"] = 3.9  # Adding GPA
        \end{lstlisting}
    \end{block}
    
    \begin{block}{Result}
    \begin{lstlisting}[language=Python]
print(student)
# Output: {'name': 'Alice', 'age': 22, 'major': 'Computer Science', 'GPA': 3.9}
    \end{lstlisting}
    \end{block}
\end{frame}

\begin{frame}[fragile]
    \frametitle{Modifying and Removing Items}
    \begin{block}{Modifying Items}
        Change the value of an existing key by direct assignment.
        \begin{lstlisting}[language=Python]
student["age"] = 23  # Updating age
        \end{lstlisting}
    \end{block}
    
    \begin{block}{Example Output}
    \begin{lstlisting}[language=Python]
print(student)
# Output: {'name': 'Alice', 'age': 23, 'major': 'Computer Science', 'GPA': 3.9}
    \end{lstlisting}
    \end{block}
\end{frame}

\begin{frame}[fragile]
    \frametitle{Removing Items}
    \begin{block}{Removing Items}
        Use \texttt{del} or the \texttt{.pop()} method to remove an item.
    \end{block}
    
    \begin{block}{Examples}
        Using \texttt{del}:
        \begin{lstlisting}[language=Python]
del student["major"]  # Removes the major key
        \end{lstlisting}
        
        Using \texttt{.pop()}:
        \begin{lstlisting}[language=Python]
gpa = student.pop("GPA")  # Removes GPA and returns its value
        \end{lstlisting}
    \end{block}
\end{frame}

\begin{frame}[fragile]
    \frametitle{Accessing Values}
    \begin{block}{Accessing Items}
        Access values using keys inside square brackets \texttt{[]} or using \texttt{.get()}.
    \end{block}
    
    \begin{block}{Example}
    \begin{lstlisting}[language=Python]
name = student["name"]        # Accessing name
age = student.get("age")      # Accessing age
    \end{lstlisting}
    \end{block}
\end{frame}

\begin{frame}[fragile]
    \frametitle{Key Points and Summary}
    \begin{block}{Key Points to Emphasize}
        \begin{itemize}
            \item Keys must be unique: Adding a duplicate key will overwrite the existing value.
            \item Dictionaries are unordered until Python 3.7; they maintain insertion order afterward.
            \item Accessing non-existent keys: Using \texttt{[]} raises a KeyError; \texttt{get()} allows specifying a default value.
        \end{itemize}
    \end{block}
    
    \begin{block}{Quick Summary}
        \begin{itemize}
            \item Creating: \texttt{my\_dict = \{key: value\}}
            \item Adding: \texttt{my\_dict[new\_key] = new\_value}
            \item Modifying: \texttt{my\_dict[existing\_key] = new\_value}
            \item Removing: \texttt{del my\_dict[key]} or \texttt{my\_dict.pop(key)}
            \item Accessing: \texttt{value = my\_dict[key]} or \texttt{value = my\_dict.get(key, default)}
        \end{itemize}
    \end{block}
\end{frame}

\begin{frame}[fragile]
    \frametitle{Common Dictionary Methods - Overview}
    Dictionaries in Python are powerful data structures that store key-value pairs. 
    Understanding the built-in methods for dictionaries is crucial for effective data manipulation and retrieval. 
    In this slide, we will cover the following essential dictionary methods:
    \begin{itemize}
        \item \texttt{.get()}
        \item \texttt{.keys()}
        \item \texttt{.values()}
        \item \texttt{.items()}
    \end{itemize}
\end{frame}

\begin{frame}[fragile]
    \frametitle{Common Dictionary Methods - The \texttt{.get()} Method}
    \begin{block}{Purpose}
        Safely retrieves a value for a given key. If the key does not exist, it returns \texttt{None} or a specified default value instead of raising an error.
    \end{block}
    \textbf{Syntax:} \texttt{dictionary.get(key, default\_value)}

    \begin{block}{Example}
    \begin{lstlisting}[language=Python]
sample_dict = {'a': 1, 'b': 2}
print(sample_dict.get('b'))        # Output: 2
print(sample_dict.get('c', 'Not Found'))  # Output: Not Found
    \end{lstlisting}
    \end{block}
\end{frame}

\begin{frame}[fragile]
    \frametitle{Common Dictionary Methods - The \texttt{.keys()}, \texttt{.values()}, and \texttt{.items()} Methods}
    \begin{block}{The \texttt{.keys()} Method}
        Returns a view object that displays a list of all the keys in the dictionary.
        \textbf{Syntax:} \texttt{dictionary.keys()}
        
        \textbf{Example:}
        \begin{lstlisting}[language=Python]
person = {'name': 'Alice', 'age': 25, 'city': 'New York'}
key_list = person.keys()
print(key_list)  # Output: dict_keys(['name', 'age', 'city'])
        \end{lstlisting}
    \end{block}

    \begin{block}{The \texttt{.values()} Method}
        Returns a view object that displays a list of all the values in the dictionary.
        \textbf{Syntax:} \texttt{dictionary.values()}
        
        \textbf{Example:}
        \begin{lstlisting}[language=Python]
value_list = person.values()
print(value_list)  # Output: dict_values(['Alice', 25, 'New York'])
        \end{lstlisting}
    \end{block}

    \begin{block}{The \texttt{.items()} Method}
        Returns a view object that displays a list of the dictionary’s key-value pairs as tuples.
        \textbf{Syntax:} \texttt{dictionary.items()}
        
        \textbf{Example:}
        \begin{lstlisting}[language=Python]
items_list = person.items()
print(items_list)  # Output: dict_items([('name', 'Alice'), ('age', 25), ('city', 'New York')])
        \end{lstlisting}
    \end{block}
\end{frame}

\begin{frame}
    \frametitle{Introduction to Sets}
    \textbf{Definition:} A set is an unordered collection of unique, immutable items in Python. 
    \begin{itemize}
        \item Allows efficient membership testing
        \item Eliminates duplicate entries
        \item Supports mathematical set operations: union, intersection, difference
    \end{itemize}
\end{frame}

\begin{frame}
    \frametitle{Unique Properties of Sets}
    \begin{enumerate}
        \item \textbf{Unordered:} Items do not have defined order; inaccessible by index
        \item \textbf{Unique Elements:} Automatically discards duplicates
        \item \textbf{Mutable:} Can modify the set but elements must be immutable
        \item \textbf{Dynamic Size:} Grows and shrinks with additions or removals
    \end{enumerate}
\end{frame}

\begin{frame}
    \frametitle{Sets vs. Lists and Dictionaries}
    \textbf{Lists:}
    \begin{itemize}
        \item Ordered collection; allows duplicates
        \item Accessed via index
        \item \texttt{my\_list = [1, 2, 2, 3]}
    \end{itemize}

    \textbf{Dictionaries:}
    \begin{itemize}
        \item Collection of key-value pairs; keys are unique
        \item Unordered; fast retrieval using keys
        \item \texttt{my\_dict = \{'a': 1, 'b': 2\}}
    \end{itemize}
\end{frame}

\begin{frame}[fragile]
    \frametitle{Code Snippet Demonstrating Sets}
    \begin{lstlisting}[language=Python]
# Creating a set with unique elements
my_set = {1, 2, 3, 4, 5}
print(my_set)  # Output: {1, 2, 3, 4, 5}

# Adding an element
my_set.add(6)
print(my_set)  # Output: {1, 2, 3, 4, 5, 6}

# Attempting to add a duplicate element
my_set.add(3)
print(my_set)  # Output: {1, 2, 3, 4, 5, 6} (No duplicates)
    \end{lstlisting}
\end{frame}

\begin{frame}
    \frametitle{Key Points}
    \begin{itemize}
        \item Best used for collections of unique items
        \item Efficient for membership testing
        \item Set operations simplify complex tasks involving multiple collections
    \end{itemize}
\end{frame}

\begin{frame}
    \frametitle{Conclusion}
    Understanding sets is crucial in Python programming. They simplify tasks related to collections and significantly contribute to effective data manipulation.
\end{frame}

\begin{frame}[fragile]
    \frametitle{Creating and Using Sets - What is a Set?}
    \begin{itemize}
        \item A \textbf{set} is a built-in data structure in Python that holds an \textbf{unordered} collection of \textbf{unique} elements.
        \item Uniqueness means that no two items in a set can be the same.
        \item Sets are mutable, allowing changes to their contents after creation, but they do not support indexing.
    \end{itemize}
\end{frame}

\begin{frame}[fragile]
    \frametitle{Creating Sets}
    \begin{block}{Using Curly Braces}
        A set can be created by placing comma-separated values inside curly braces \{\}:
        \begin{lstlisting}[language=Python]
my_set = {1, 2, 3}
        \end{lstlisting}
    \end{block}

    \begin{block}{Using the set() Constructor}
        Create a set from a list or other iterable using the \texttt{set()} function:
        \begin{lstlisting}[language=Python]
my_list = [1, 2, 2, 3]
my_set = set(my_list)  # Result: {1, 2, 3}
        \end{lstlisting}
    \end{block}
\end{frame}

\begin{frame}[fragile]
    \frametitle{Adding and Removing Elements}
    \begin{itemize}
        \item \textbf{Adding Elements:}
        \begin{itemize}
            \item Use \texttt{.add()} to add a single element:
            \begin{lstlisting}[language=Python]
my_set.add(4)  # Now my_set is {1, 2, 3, 4}
            \end{lstlisting}
            \item Use \texttt{.update()} to add multiple elements:
            \begin{lstlisting}[language=Python]
my_set.update([5, 6])  # Now my_set is {1, 2, 3, 4, 5, 6}
            \end{lstlisting}
        \end{itemize}

        \item \textbf{Removing Elements:}
        \begin{itemize}
            \item Use \texttt{.remove()}:
            \begin{lstlisting}[language=Python]
my_set.remove(2)  # Now my_set is {1, 3, 4, 5, 6}
            \end{lstlisting}
            \item Use \texttt{.discard()} (does not raise an error if not found):
            \begin{lstlisting}[language=Python]
my_set.discard(10)  # No error, my_set remains unchanged
            \end{lstlisting}
            \item Use \texttt{.pop()} (removes an arbitrary element):
            \begin{lstlisting}[language=Python]
popped_element = my_set.pop()  # Removes and returns an element
            \end{lstlisting}
        \end{itemize}
    \end{itemize}
\end{frame}

\begin{frame}[fragile]
    \frametitle{Set Operations and Key Points}
    \begin{itemize}
        \item \textbf{Set Operations:}
        \begin{itemize}
            \item \textbf{Union} (\(A \cup B\)): Combines unique elements from both sets (\(A | B\)).
            \item \textbf{Intersection} (\(A \cap B\)): Elements present in both sets (\(A & B\)).
            \item \textbf{Difference} (\(A - B\)): Elements in set \(A\) that are not in set \(B\).
        \end{itemize}
        \item \textbf{Key Points to Remember:}
        \begin{itemize}
            \item Sets are unordered and do not allow duplicates.
            \item Elements can be added or removed dynamically.
            \item Essential for mathematical set operations in programming.
            \item Useful in scenarios where the presence of distinct elements is crucial (e.g., collecting unique user IDs).
        \end{itemize}
    \end{itemize}
\end{frame}

\begin{frame}[fragile]
    \frametitle{Next Topic}
    \begin{itemize}
        \item Review common set methods including \texttt{.union()}, \texttt{.intersection()}, and \texttt{.difference()}.
    \end{itemize}
\end{frame}

\begin{frame}
    \frametitle{Common Set Methods}
    \begin{block}{Overview}
        Sets are a powerful data structure in Python that allow for the storage of unique elements. Here, we will explore three common set methods: 
        \begin{itemize}
            \item \texttt{.union()}
            \item \texttt{.intersection()}
            \item \texttt{.difference()}
        \end{itemize}
        Understanding these methods will enhance your ability to perform operations on sets effectively.
    \end{block}
\end{frame}

\begin{frame}[fragile]
    \frametitle{Union (\texttt{.union()})}
    \begin{block}{Definition}
        Combines all unique elements from two or more sets. If there are duplicate elements, they are included only once in the resulting set.
    \end{block}
    \begin{block}{Syntax}
        \texttt{set1.union(set2) OR set1 \textbar{} set2}
    \end{block}
    
    \begin{block}{Example}
    \begin{lstlisting}[language=Python]
set_a = {1, 2, 3}
set_b = {3, 4, 5}
result = set_a.union(set_b)
# Output: {1, 2, 3, 4, 5}
    \end{lstlisting}
    \end{block}
\end{frame}

\begin{frame}[fragile]
    \frametitle{Intersection (\texttt{.intersection()}) and Difference (\texttt{.difference()})}
    \begin{block}{Intersection (\texttt{.intersection()})}
        \begin{itemize}
            \item \textbf{Definition:} Returns only the elements that are common to both sets.
            \item \textbf{Syntax:} \texttt{set1.intersection(set2) OR set1 \& set2}
        \end{itemize}
        \begin{block}{Example}
        \begin{lstlisting}[language=Python]
set_a = {1, 2, 3}
set_b = {3, 4, 5}
result = set_a.intersection(set_b)
# Output: {3}
        \end{lstlisting}
        \end{block}
    \end{block}

    \begin{block}{Difference (\texttt{.difference()})}
        \begin{itemize}
            \item \textbf{Definition:} Returns the elements that are in the first set but not in the second set.
            \item \textbf{Syntax:} \texttt{set1.difference(set2) OR set1 - set2}
        \end{itemize}
        \begin{block}{Example}
        \begin{lstlisting}[language=Python]
set_a = {1, 2, 3}
set_b = {3, 4, 5}
result = set_a.difference(set_b)
# Output: {1, 2}
        \end{lstlisting}
        \end{block}
    \end{block}
\end{frame}

\begin{frame}
    \frametitle{Key Points and Conclusion}
    \begin{block}{Key Points}
        \begin{itemize}
            \item \textbf{Uniqueness:} Sets only store unique elements, making them ideal for membership testing and removing duplicates.
            \item \textbf{Mutability:} Sets are mutable, meaning you can change their contents (add/remove elements) after they are created.
            \item \textbf{Operations:} Use these methods to combine, find commonalities, or differentiate datasets efficiently.
        \end{itemize}
    \end{block}
    \begin{block}{Conclusion}
        Utilizing these common set methods expands your ability to manipulate collections of unique items effectively in Python, facilitating a wide range of applications in data analysis and processing.
    \end{block}
\end{frame}

\begin{frame}
    \frametitle{Practical Applications in Data Manipulation}
    \begin{block}{Introduction to Dictionaries and Sets}
        \begin{itemize}
            \item \textbf{Dictionaries}: Data structure storing data in key-value pairs. Ideal for fast lookup, insertion, and deletion using unique keys.
            \item \textbf{Sets}: Unordered collection of unique elements. Useful for eliminating duplicates and performing set operations.
        \end{itemize}
    \end{block}
\end{frame}

\begin{frame}[fragile]
    \frametitle{Real-World Applications - Part 1}
    \begin{enumerate}
        \item \textbf{Dictionaries in User Profiles}
            \begin{lstlisting}[language=Python]
user_profile = {
    "username": "johndoe",
    "email": "johndoe@example.com",
    "age": 30,
    "preferences": ["news", "sports", "music"]
}
            \end{lstlisting}
            \item \textbf{Use Case}: Easily retrieve or update user information (e.g., fetching a user’s email based on their username).
        
        \item \textbf{Sets for Unique Item Tracking}
            \begin{lstlisting}[language=Python]
shopping_cart = {"apple", "banana", "orange"}
            \end{lstlisting}
            \item \textbf{Use Case}: Prevent adding duplicate items in the shopping cart. E.g., trying to add "apple" again remains unchanged.
    \end{enumerate}
\end{frame}

\begin{frame}[fragile]
    \frametitle{Real-World Applications - Part 2}
    \begin{enumerate}
        \setcounter{enumi}{2} % Resume enumeration from 3
        \item \textbf{Dictionaries for Frequency Counting}
            \begin{lstlisting}[language=Python]
text = "hello world hello"
word_count = {}
for word in text.split():
    word_count[word] = word_count.get(word, 0) + 1
            \end{lstlisting}
            \item \textbf{Use Case}: Quick access to how many times each word appears for text analysis.
        
        \item \textbf{Sets for Data Deduplication}
            \begin{lstlisting}[language=Python]
email_list = ["test@example.com", "hello@example.com", "test@example.com"]
unique_emails = set(email_list)  # {'test@example.com', 'hello@example.com'}
            \end{lstlisting}
            \item \textbf{Use Case}: Filters out duplicate entries to ensure valid email addresses are stored.
    \end{enumerate}
\end{frame}

\begin{frame}
    \frametitle{Key Points and Conclusion}
    \begin{block}{Key Points to Emphasize}
        \begin{itemize}
            \item \textbf{Efficiency}: O(1) average time complexity for lookups in dictionaries and sets.
            \item \textbf{Flexibility}: Adaptable to various applications in finance, e-commerce, and user management.
            \item \textbf{Data Integrity}: Maintains data uniqueness and facilitates quick modifications.
        \end{itemize}
    \end{block}
    
    \begin{block}{Conclusion}
        Dictionaries and sets are essential tools for organizing and manipulating data, enhancing performance and data management capabilities for developers and data scientists alike.
    \end{block}
\end{frame}

\begin{frame}
    \frametitle{Group Project Overview}
    \begin{block}{Objective}
        Utilize dictionaries and sets in Python to collaboratively solve a real-world problem.
    \end{block}
\end{frame}

\begin{frame}
    \frametitle{Project Guidelines}
    \begin{enumerate}
        \item \textbf{Project Selection:}
        \begin{itemize}
            \item Choose a relevant problem that could benefit from data organization or manipulation. 
            \item Examples include:
                \begin{itemize}
                    \item Analyzing survey data
                    \item Managing inventory for a small business
                    \item Tracking student grades
                    \item Collecting and organizing environmental data
                \end{itemize}
        \end{itemize}
        
        \item \textbf{Utilization of Data Structures:}
        \begin{itemize}
            \item \textbf{Dictionaries:} Use for key-value pairs to map relationships. 
            \item \textbf{Sets:} Use for collections of unique items, useful for various tasks.
        \end{itemize}
    \end{enumerate}
\end{frame}

\begin{frame}
    \frametitle{Team Roles and Deliverables}
    \begin{block}{Team Roles}
        \begin{itemize}
            \item \textbf{Project Manager:} Oversees project timeline and deliverables.
            \item \textbf{Data Analyst:} Focuses on data manipulation using dictionaries and sets.
            \item \textbf{Moderator:} Manages discussions and synthesizes ideas.
            \item \textbf{Presenter:} Develops and delivers the final presentation.
        \end{itemize}
    \end{block}

    \begin{block}{Deliverables}
        \begin{itemize}
            \item \textbf{Final Report (2-3 pages):} Summarize your project findings, methodology, and outcomes.
            \item \textbf{Code Implementation:} A working script demonstrating the functionality of your solution.
            \item \textbf{Presentation (10-15 minutes):} Present project details and findings, emphasizing your use of dictionaries and sets.
        \end{itemize}
    \end{block}
\end{frame}

\begin{frame}[fragile]
    \frametitle{Example Code Snippet}
    \begin{lstlisting}[language=Python]
# Using a dictionary to store student grades
grades = {
    "Alice": [90, 85, 92],
    "Bob": [78, 82, 88],
    "Charlie": [85, 90, 91]
}

# Using a set to find unique subjects taken
subjects = set(["Math", "Science", "English", "Art", "Math"])
print(subjects)  # Output: {'Math', 'Science', 'English', 'Art'}
    \end{lstlisting}
\end{frame}

\begin{frame}
    \frametitle{Conclusion and Next Steps}
    \begin{block}{Conclusion}
        This group project will deepen your understanding of dictionaries and sets through practical application. Embrace creativity and collaboration to find innovative solutions to your chosen problem!
    \end{block}

    \begin{block}{Next Steps}
        Review Chapter 6 for deeper insights into dictionaries and sets, and prepare ideas for your project topic before the next class.
    \end{block}
\end{frame}

\begin{frame}[fragile]
    \frametitle{Summary and Q\&A}
    \begin{block}{Recap: Key Concepts of Dictionaries and Sets}
        \begin{itemize}
            \item \textbf{Dictionaries:}
                \begin{itemize}
                    \item A collection of key-value pairs where each key is unique.
                    \item Created using curly braces \{\} or the \texttt{dict()} function.
                    \item Access values using keys.
                    \item Common methods: 
                        \begin{itemize}
                            \item \texttt{my\_dict.keys()} – Returns all keys.
                            \item \texttt{my\_dict.values()} – Returns all values.
                            \item \texttt{my\_dict.items()} – Returns key-value pairs.
                        \end{itemize}
                \end{itemize}
            \item \textbf{Sets:}
                \begin{itemize}
                    \item An unordered collection of unique elements.
                    \item Created using curly braces \{\} or the \texttt{set()} function.
                    \item Operations include union, intersection, and difference.
                \end{itemize}
        \end{itemize}
    \end{block}
\end{frame}

\begin{frame}[fragile]
    \frametitle{Key Points and Applications}
    \begin{block}{Key Points to Emphasize}
        \begin{itemize}
            \item \textbf{Uniqueness:} 
                \begin{itemize}
                    \item Dictionaries (keys) and sets (elements) do not allow duplicates. 
                \end{itemize}
            \item \textbf{Efficiency:}
                \begin{itemize}
                    \item Efficient operations for insertion, deletion, and membership tests.
                \end{itemize}
            \item \textbf{Applications:} 
                \begin{itemize}
                    \item Data modeling, caching results, and performing statistical operations on datasets.
                \end{itemize}
        \end{itemize}
    \end{block}
\end{frame}

\begin{frame}[fragile]
    \frametitle{Q\&A Session}
    \begin{block}{Encourage Questions}
        \begin{itemize}
            \item Invite students to ask questions regarding:
                \begin{itemize}
                    \item Implementation of dictionaries and sets in Python.
                    \item Differences between dictionaries and sets.
                    \item Scenarios for use in development.
                \end{itemize}
            \item \textbf{Discussion Points:}
                \begin{itemize}
                    \item When to use a dictionary vs. a set.
                    \item Examples of real applications in software development.
                \end{itemize}
        \end{itemize}
    \end{block}
    
    \begin{block}{Example Code Snippet}
        \begin{lstlisting}[language=Python]
# Merging dictionaries
dict_a = {'a': 1, 'b': 2}
dict_b = {'b': 3, 'c': 4}
merged_dict = {**dict_a, **dict_b}
print(merged_dict)  # Output: {'a': 1, 'b': 3, 'c': 4}
        \end{lstlisting}
    \end{block}
\end{frame}


\end{document}